\section{The strict mixed-fibre component of the original phase observation}
\label{GraphV2:sec:pgm-strict-mixed}

Retain every object and coordinate of PGS.1--PGS.41. In particular
\(d=4m\), \(m\geq1\), \(k\geq1\), \(a=kA>0\),
\[
 s_i=\tfrac12+it_i,\qquad
 u=\sum_{i=1}^kt_i,\qquad Q(t)=\sum_{i=1}^kq_\Phi(t_i),\qquad
 \tau_k=a+iQ(t),\qquad c=1+a^2 .
 \tag{PGM.1}\label{GraphV2:eq:pgm-data}
\]
The real odd polynomial \(q_\Phi\) has degree \(d+1\) and leading
coefficient \(1/(d+1)\). The full original cyclic source is
\(E=\mathbb{C}[S]/(\chi_{h,k})\) of dimension
\(q=[1+k(m-1)](k+1)^2\). Its actual observation is
\[
\begin{gathered}
 Lx=\Bigl(\prod_i v_h(s_i)\Bigr)F_x(s_1,\ldots,s_k),\qquad
 \\
 F_x=\sum_{b=0}^{q-1}x_b
       \operatorname{rem}_{h(s_1),\ldots,h(s_k)}
             \left(\Bigl[\sum_i\Phi(s_i)\Bigr]\Bigl[\sum_i s_i\Bigr]^b\right).
\end{gathered}
 \tag{PGM.2}\label{GraphV2:eq:pgm-F}
\]
Every variable degree of \(F_x\) is at most \(d-1\). The full
unit \(v_h=2\xi/h\) is retained in \(L\), which is injective by
PGS.4 and its proof: the original cyclic map and section are injective
and the actual ordered phase sums are nonzero. Thus \(F_x=0\)
implies \(x=0\). No tensor tuple, repeated critical value or
arithmetic nilpotent has been removed.

Write \(w(t)=|v_h(1/2+it)|^2/(2\pi)\). Define the two original
weighted Hilbert spaces and the sum-coordinate inclusion by
\[
\begin{gathered}
 \mathscr K_k=L^2\left(\mathbb{R}^k,
                   (c+Q(t)^2)\prod_i w(t_i)\,dt_i\right),\qquad
 \mathscr L_k=L^2(\mathbb{R},r_k(u)\,du),\\
 \jmath_k:\mathscr L_k\longrightarrow\mathscr K_k,\qquad
          (\jmath_k f)(t)=f\left(\sum_i t_i\right).
\end{gathered}
\tag{PGM.3}\label{GraphV2:eq:pgm-hilbert}
\]
The fibre density \(r_k\) is exactly PGS.35, including all three
convolutions and the original measure factors. Its fibre integral
proves \(\|\jmath_k f\|_{\mathscr K_k}=\|f\|_{\mathscr L_k}\),
so \(\jmath_k\) has closed range and
\(\mathsf P_k=\jmath_k\jmath_k^*\) is its orthogonal projection.
Define the full multiplication-graph minimizer and its mixed part:
\[
\begin{gathered}
 B_k:E\longrightarrow\mathscr K_k,\qquad
 B_kx=\frac{(a-iQ)F_x(s)}{c+Q^2},\qquad
 g_x(u)=\frac{z(u)x}{r_k(u)},\\
 \jmath_k^*B_kx=g_x,\qquad
 \mathscr D_k=(I-\mathsf P_k)B_k,\qquad
 \mathscr D_kx=B_kx-\jmath_k g_x .
\end{gathered}
\tag{PGM.4}\label{GraphV2:eq:pgm-map}
\]
Here \(B_k\) is a new map with the explicitly displayed domain and
target, not the relation-column map \(B_N\) of PGS.3. It is
well-defined, since
\[
 \|B_kx\|_{\mathscr K_k}^2
 =\int\frac{a^2+Q^2}{c+Q^2}|F_x(s)|^2\prod_iw(t_i)\,dt_i
 \leq\|Lx\|_{\mathscr H_k}^2.
\]
Pairing \(B_kx\) against \(\jmath_k f\) cancels exactly the
positive factor \(c+Q^2\); the resulting fibre integral is
\(\int\overline{f(u)}z(u)x\,du\). This proves the adjoint
formula in PGM.4, with its conjugate phase and its original masses.
The cancellation does not change either source vector.

The two residuals of PGS.17 and PGS.40 satisfy the exact matrix
identity
\[
 \boxed{\quad \mathscr D_k^*\mathscr D_k
                       =\Omega_k-K_{\rm low}.\quad}
 \tag{PGM.5}\label{GraphV2:eq:pgm-gram}
\]
Indeed \(B_k^*B_k=K-K_{\rm low}\) by the preceding integral,
where \(K=L^*L\) is the actual cyclic observation Gram.
The adjoint formula gives
\((\jmath_k^*B_k)^*(\jmath_k^*B_k)=
\int z(u)^*z(u)/r_k(u)\,du=K-\Omega_k\).
Subtracting these equalities proves PGM.5, with all cross entries
retained. It also proves the norm identity in PGS.41 directly.

\subsection{The exact kernel for every mixed tensor}

For \(k\geq2\) the map \(\mathscr D_k\) is injective. We prove this
on its actual measurable source, not on an auxiliary polynomial
model. Suppose \(\mathscr D_kx=0\). The nonzero entire function
\(v_h\) has only isolated zeros on the original line. Hence
\(\prod_iw(t_i)>0\) for Lebesgue-almost every \(t\in\mathbb{R}^k\),
and \(c+Q^2\) is everywhere strictly positive. The asserted
Hilbert-space equality therefore gives
\[
 (a-iQ(t))F_x(s)=g_x(u)(c+Q(t)^2)
                  \quad\text{for almost every }t\in\mathbb{R}^k.
 \tag{PGM.6}\label{GraphV2:eq:pgm-fibre-identity}
\]
The determinant-one change of variables
\((t_1,\ldots,t_k)\mapsto(t_1,\ldots,t_{k-1},u)\)
preserves null sets. By Fubini, for almost every \(u\) at which
the finite representative \(g_x(u)\) is defined, PGM.6 holds for
almost every \((t_1,\ldots,t_{k-1})\).

Fix such a \(u\) and, when \(k>2\), fix almost every
\((t_2,\ldots,t_{k-1})\) to which Fubini applies. Put
\(b=u-\sum_{i=2}^{k-1}t_i\), where the empty sum at \(k=2\)
is zero, and set \(t_k=b-t_1\). Exclude \(b=0\); this removes
only a null set, since it is \(u=0\) for \(k=2\), and a proper
affine hyperplane in the remaining variables for \(k>2\).
The polynomial
\[
 Q_b(X)=q_\Phi(X)+q_\Phi(b-X)
                            +\sum_{i=2}^{k-1}q_\Phi(t_i)
 \tag{PGM.7}\label{GraphV2:eq:pgm-fibre-degree}
\]
has degree exactly \(d\) and leading coefficient \(b\ne0\).
To verify this, \(d\) is even and the leading phase terms give
\([X^{d+1}+(b-X)^{d+1}]/(d+1)
=bX^d+\text{terms of lower degree}\).
Every lower term of the odd phase has degree at most \(d-1\);
its substitution cannot supply an \(X^d\) coefficient.
Meanwhile the substituted polynomial \(F_b(X)\), obtained from
the very same \(F_x\) in PGM.2, has degree at most \(2d-2\):
only its first and last variables depend on \(X\), each with
degree at most \(d-1\).

The almost-everywhere equality in \(X\) is an identity of complex
polynomials, because a nonzero one-variable polynomial has only
finitely many real roots. Set \(g=g_x(u)\in\mathbb{C}\).
Its exact polynomial consequence is
\[
\begin{gathered}
 (a-iQ_b)F_b=g(c+Q_b^2),\\
 (c+Q_b^2)-(a+iQ_b)(a-iQ_b)=1,\\
 F_b=(c+Q_b^2)\,[F_b-g(a+iQ_b)].
\end{gathered}
\tag{PGM.8}\label{GraphV2:eq:pgm-divisibility}
\]
The second identity uses the original constant \(c=1+a^2\);
the \(1\) is essential. The third follows by multiplying the
second by \(F_b\) and using the first. The polynomial
\(c+Q_b^2\) has degree \(2d\) and nonzero leading coefficient
\(b^2\). If the bracket were nonzero, its product would have
degree at least \(2d\), exceeding \(2d-2\). Therefore the
bracket and \(F_b\) are both zero.

It follows for almost every fixed fibre and remaining coordinate
choice that \(F_x(s)=0\) for every \(t_1\). Fubini and the same
invertible linear change of variables give \(F_x(s)=0\) for
almost every \(t\in\mathbb{R}^k\). A polynomial vanishing almost
everywhere on \(\mathbb{R}^k\) is identically zero: apply the
one-variable fact iteratively to its coefficient polynomials,
using Fubini. The invertible affine substitutions
\(s_i=1/2+it_i\) then prove \(F_x=0\) in its original variables.
Injectivity of \(L\) proves \(x=0\), completing the kernel proof.

For \(k=1\), the sum-coordinate inclusion in PGM.3 is the
identity between the identical weighted spaces
\(r_1(t)dt=(1+A^2+q_\Phi(t)^2)w(t)dt\).
Consequently \(\mathscr D_1=0\). Combining this with PGM.5
and the preceding injectivity gives the complete assertion
\[
\begin{gathered}
 \ker\mathscr D_k=
 \begin{cases}E,&k=1,\\ \{0\},&k\geq2,\end{cases}
 \qquad
 \operatorname{rank}\mathscr D_k=q\quad(k\geq2),\\
 \Omega_1=K_{\rm low},\qquad
 \Omega_k-K_{\rm low}>0\quad(k\geq2),\qquad
 0<K_{\rm low}<\Omega_k\preceq\Delta_N
          \quad(k\geq2,\ N\geq q-1).
\end{gathered}
\tag{PGM.9}\label{GraphV2:eq:pgm-strict}
\]
Strict matrix inequalities use the original coefficient frame of
\(E\); for every nonzero \(x\), their value is the strictly
positive squared norm of the actual \(\mathscr D_kx\).

On each retained coefficient mask \(A_{\rm mask}\), the same map
acts coordinatewise from \(E^{A_{\rm mask}}\) into
\(\mathscr K_k^{A_{\rm mask}}\). Its kernel is zero for \(k\geq2\).
It commutes with zero insertion between masks because it is linear.
The present empty tuple retains its outer support label; absence
\(\tau_{\rm out}\) remains a separate point and maps to absence.
For the already specified compatible coefficient maps
\(\alpha:R\to\mathbb{C}\), \(\lambda:R\to R'\), and
\(\alpha':R'\to\mathbb{C}\) with \(\alpha'\lambda=\alpha\), set
\[
\begin{gathered}
 \mu_E:R'\otimes_R E\longrightarrow E,\qquad
       b\otimes x\longmapsto\alpha'(b)x,\\
 \mathscr D_k^{\rm real}=\mathscr D_k\mu_E,\qquad
 \mathscr D_k^{\rm real}(b\otimes x)
                         =\alpha'(b)\mathscr D_kx,\\
 \ker\mathscr D_k^{\rm real}=\ker\mu_E
   =\left\{\sum_j(b_j\otimes x_j
                    -1\otimes\alpha'(b_j)x_j)\right\}\quad(k\geq2).
\end{gathered}
\tag{PGM.10}\label{GraphV2:eq:pgm-realization}
\]
Balancedness follows from the specified equality of coefficient
actions. Every generator in the last display has zero image under
\(\mu_E\). Conversely, for a tensor \(z=\sum_jb_j\otimes x_j\),
subtracting \(1\otimes\mu_E(z)\) gives exactly the displayed sum.
The map \(x\mapsto1\otimes x\) is an \(R\)-linear section, so this
also proves surjectivity and identifies every fibre as
\(1\otimes x+\ker\mu_E\). Injectivity of \(\mathscr D_k\)
proves the equality of kernels. This constructs the mixed component
on the existing carrier, with the full supported-zero realization
fibre retained. It does not assert an unprovided residue-field map
to \(\mathbb{C}\).

The original polynomial and theta relation arrows have an equally
explicit comparison. For \(k\geq2\) and \(N\geq q-1\), put
\(\mathscr D_{k,N}=\mathscr D_kJ_N\). Then
\[
 0\longrightarrow\mathcal P_{N-q}
  \xrightarrow{\ B_N\ }\mathcal P_N
  \xrightarrow{\ \mathscr D_{k,N}\ }\mathscr D_k(E)
  \longrightarrow0,\qquad
 \mathscr D_{k,N}s_N=\mathscr D_k
 \tag{PGM.11}\label{GraphV2:eq:pgm-source-exact}
\]
is exact. Indeed \(B_NQ=\chi Q\) is injective by monic
division, \(J_N\) is onto with kernel \(\chi\mathcal P_{N-q}\),
and \(\mathscr D_k\) is injective. These statements prove
each exactness assertion and the section equality, including
\(\mathcal P_{-1}=0\) at \(N=q-1\).
On the literal Hilbert source \(\mathcal V_N(\mathcal P_N)\),
the same arrow is
\(\mathscr D_{k,N}\mathcal V_N^{-1}\), where the inverse is
defined on that actual image by the proved injectivity of
\(\mathcal V_N\). Its kernel is precisely
\(\mathcal V_N(\chi\mathcal P_{N-q})\).
Thus all nonzero original relation vectors, their theta primitives
and their original norms remain present. The arrow factors through
the stated polynomial cyclic quotient, not through an unidentified
proper-source quotient.

Finally, for \(k\geq2\), its inverse on the mixed image is explicitly
\[
\begin{gathered}
 \mathscr R_k^{\rm mix}:\mathscr K_k\longrightarrow E,\qquad
 \mathscr R_k^{\rm mix}=(\Omega_k-K_{\rm low})^{-1}\mathscr D_k^*,
 \\
 \mathscr R_k^{\rm mix}\mathscr D_k=I_E,\qquad
 \mathscr D_k\mathscr R_k^{\rm mix}
       =\text{orthogonal projection onto }\mathscr D_k(E).
\end{gathered}
 \tag{PGM.12}\label{GraphV2:eq:pgm-left-inverse}
\]
The first composite follows from PGM.5 and PGM.9. The second
operator is self-adjoint, its square equals itself, and it acts
identically on \(\mathscr D_k(E)\); its range is contained in that
image. These facts prove the projection assertion with the
unchanged weighted Hilbert metric. All adjoints and matrix inverses
refer to the original cyclic coefficient frame specified above.

The result is a strict, explicitly calculated mixed-source residual.
Its matrix entries are the integrals in PGM.4--PGM.5 with their
unchanged units and domains. It does not give a uniform estimate
for its eigenvalues as \(k\) grows. In particular it does not replace
the signed four-window identity PGS.34 by an asserted sign or by a
shrinking exponent error.
